\begin{landscape}
\begingroup
\fontsize{10}{13}\selectfont
\fancyhead[L]{\small NACHLASS. STÖRUNGEN DER PALLAS}
\renewcommand{\headrulewidth}{0pt}
\par\medskip
\begingroup\makeatletter\tagsleft@true\makeatother
\pdfbookmark[2]{Nachlass VI [11]}{pallas-VI-fragment11-first}
\begin{center}[11.]\end{center}
\makeatletter
\providecommand{\VIelevenFirstHistoricalNote}[1]{\begingroup
\renewcommand{\thefootnote}{*)}
\long\def\@makefntext##1{\noindent [*)\ ##1}
\footnote{#1}\endgroup}
\makeatother
\par\noindent\hspace*{\parindent}[Zur Vergleichung der Resultate der zweiten Rechnung der generellen Störungen mit den Oppositionen wurde im allgemeinen ebenso verfahren, wie bei den frühern Vergleichungen; es wurden aber die Argumente \(\text{\mdseries\AstroSymbols ♃}\) und \(\text{\mdseries\AstroSymbols ⚴}\) für die Störungsgleichungen schärfer, nemlich mit Berücksichtigung der grossen Gleichung für beide Planeten berechnet nach den Formeln
\[
\begin{aligned}
\text{\mdseries\AstroSymbols ⚴}&=47^0 24'\,28\overset{\prime\prime}{,}37+769\overset{\prime\prime}{,}15194t+\text{Grosse Gleichung}\\
\text{\mdseries\AstroSymbols ♃}&=25\quad24\quad44{,}12+299{,}12817t+\text{Grosse Gleichung,}
\end{aligned}
\tag*{1)}
\]
wo die Zeit \(t\) von 1810 Januar 0. Mittl. Zeit Göttingen gezählt ist. Die grosse Gleichung für Jupiter wurde aus den \textsc{Bouvard}schen Tafeln entnommen: die der Pallas aber in folgender Weise ermittelt.
\par\noindent\hspace*{\parindent}Setzt man
\[
u=5\text{\mdseries\AstroSymbols ♃}-2\text{\mdseries\AstroSymbols ⚴},
\]
so ist der aus \(\displaystyle\int n\mathrm{d}t\) stammende Theil der grossen Gleichung
\[
\begin{aligned}
\mathrm{f}(u)={}&(3{,}49256)\sin(u+75^0 35'\,33'')\\
&+(1{,}95199)\sin(2u+145\quad56\quad15\phantom{'})\\
&+(0{,}65379)\sin(3u+216\quad46\quad43\phantom{'})
\end{aligned}
\]
und der aus \(\varepsilon\) stammende
\[
\begin{aligned}
\mathrm{F}(u)={}&(2{,}38474)\sin(u+73^0 14'\,52'')\\
&+(1{,}14525)\sin(2u+143\quad9\quad51\phantom{'})\\
&+(0{,}01994)\sin(3u+214\quad17\quad10\phantom{'})
\end{aligned}
\]
oder summirt die grosse Gleichung
\[
\begin{aligned}
\theta={}&(3{,}52516)\sin(u+75^0 25'\,22'')\\
&+(2{,}01491)\sin(2u+145\quad33\quad48\phantom{'})\\
&+(0{,}74446)\sin(3u+216\quad18\quad31\phantom{'}),
\end{aligned}
\tag*{2)}
\]
wo die Zahlen in Klammern Logarithmen sind.
\par\noindent\hspace*{\parindent}Bezeichnet man mit \(u_0\) den Werth, den \(u\) bei Vernachlässigung von \(\theta\) annimmt, und den man also linear aus den mittlern Bewegungen erhält, so lässt sich \(\theta\) nach Vielfachen von \(u_0\) entwickeln. Diese Entwickelung wurde nach der oft angewendeten interpolatorischen Methode ausgeführt, indem in der \SourceLeaf{566}vorstehenden Gleichung zunächst \(u_0\) für \(u\) gesetzt, und \(\theta\) für die 12 Werthe \(0^0,30^0,60^0\,.\,.\,.\,.\,.\,330^0\) von \(u_0\) berechnet wurde; hiemit fanden sich Näherungswerthe für \(u=u_0-2\theta\), mit denen die Berechnung von \(\theta\) wiederholt wurde. Die Interpolation ergab schliesslich:
\[
\begin{aligned}
\theta={}&+3252\overset{\prime\prime}{,}56\cos u_0&&+846\overset{\prime\prime}{,}85\sin u_0\\
&+\phantom{32}51{,}64\cos2u_0&&-\phantom{8}72{,}93\sin2u_0\\
&-\phantom{325}2{,}16\cos3u_0&&-\phantom{84}3{,}15\sin3u_0\\
&+\phantom{325}0{,}16\cos4u_0&&-\phantom{84}0{,}06\sin4u_0.
\end{aligned}
\tag*{3)}
\]
Zur Erleichterung der Rechnung wurde eine Hülfstafel hergestellt, welche \(\theta\) für das Argument \(u_0\) von Grad zu Grad gibt.
\par\noindent\hspace*{\parindent}Für die Argumente ergaben sich hiemit nach 1) die Werthe:
\begin{center}
\begingroup\fontsize{8}{10}\selectfont\setlength{\tabcolsep}{2pt}
\begin{tabular}{@{}lr|r|r|r|r|r|r|r@{}}
Opposition:&I&II&III&IV&V&VI&VII&VIII\\
\(\text{\mdseries\AstroSymbols ⚴}=\)&\(260^0 24'\,10''\)&\(351^0 43'\,15''\)&\(89^0 16'\,43''\)&\(200^0 40'\,50''\)&\(296^0 42'\,38''\)&\(27^0\ 4'\,15''\)&\(137^0 35'\,6''\)&\(238^0 55'\,53''\)\\
\(\text{\mdseries\AstroSymbols ♃}=\)&\(188\ 19\ 23\)&\(223\ 49\ 6\)&\(261\ 43\ 46\)&\(305\ 1\ 55\)&\(342\ 21\ 44\)&\(17\ 29\ 31\)&\(60\ 27\ 39\)&\(99\ 52\ 19\)
\end{tabular}
\endgroup
\end{center}
und hiemit die Beträge der allgemeinen Störungen:
\begin{center}
\begingroup\fontsize{8}{10}\selectfont\setlength{\tabcolsep}{3pt}
\begin{tabular}{@{}lr|r|r|r|r|r|r|r@{}}
\(\delta\AscendingNode=\)&\(+146\overset{\prime\prime}{,}94\)&\(+164\overset{\prime\prime}{,}75\)&\(+146\overset{\prime\prime}{,}61\)&\(+133\overset{\prime\prime}{,}55\)&\(+104\overset{\prime\prime}{,}15\)&\(+102\overset{\prime\prime}{,}96\)&\(-15\overset{\prime\prime}{,}12\)&\(-52\overset{\prime\prime}{,}83\)\\
\(\delta i=\)&\(+93{,}30\)&\(+81{,}59\)&\(+58{,}26\)&\(+72{,}06\)&\(+77{,}13\)&\(+24{,}20\)&\(-61{,}25\)&\(-76{,}94\)\\
\(\delta\varphi=\)&\(+583{,}52\)&\(+509{,}42\)&\(+474{,}89\)&\(+426{,}91\)&\(+308{,}92\)&\(-78{,}90\)&\(-238{,}71\)&\(-228{,}76\)\\
\(\delta\tilde\omega=\)&\(+344{,}59\)&\(+267{,}36\)&\(+61{,}89\)&\(-88{,}60\)&\(-173{,}50\)&\(-995{,}19\)&\(-844{,}85\)&\(-576{,}78.\)
\end{tabular}
\endgroup
\end{center}
\par\noindent\hspace*{\parindent}Die mittlern Elemente können nun mit Benutzung der Resultate der speciellen Störungen bestimmt werden. Aus den Normalelementen (geltend für 1810)\VIelevenFirstHistoricalNote{Diese Werthe hatte eine Vergleichung der speciellen Störungen mit den ersten 8 Oppositionen ergeben, welche nach der S. 486—488 abgedruckten ausgeführt wurde.]}:
\[
\begin{aligned}
\AscendingNode&=172^0 34'\,12\overset{\prime\prime}{,}40&\qquad\varphi&=14^0 12'\,42\overset{\prime\prime}{,}21\\
i&=34\quad37\quad43{,}09&\tilde\omega&=121\quad13\quad27{,}33
\end{aligned}
\]
folgen nach Anbringung der Störungen S. 486 die osculirenden Elemente:
\begin{center}
\begingroup\fontsize{8}{10}\selectfont\setlength{\tabcolsep}{2pt}
\begin{tabular}{@{}lr|r|r|r|r|r|r|r@{}}
\(\AscendingNode=\)&\(172^0 34'\,12\overset{\prime\prime}{,}85\)&\(34'\,30\overset{\prime\prime}{,}28\)&\(34'\,13\overset{\prime\prime}{,}21\)&\(33'\,59\overset{\prime\prime}{,}92\)&\(33'\,28\overset{\prime\prime}{,}85\)&\(33'\,26\overset{\prime\prime}{,}28\)&\(31'\,28\overset{\prime\prime}{,}51\)&\(30'\,51\overset{\prime\prime}{,}14\)\\
\(i=\)&\(34\ 37\ 43{,}04\)&\(37\ 31{,}92\)&\(37\ 8{,}51\)&\(37\ 21{,}26\)&\(37\ 26{,}83\)&\(36\ 35{,}70\)&\(35\ 10{,}97\)&\(34\ 53{,}50\)\\
\(\varphi=\)&\(14\ 12\ 38{,}60\)&\(11\ 24{,}39\)&\(10\ 49{,}82\)&\(10\ 1{,}31\)&\(8\ 1{,}62\)&\(1\ 27{,}86\)&\(13^0 58\ 48{,}12\)&\(58\ 57{,}68\)\\
\(\tilde\omega=\)&\(121\ 13\ 27{,}35\)&\(12\ 11{,}96\)&\(8\ 51{,}38\)&\(6\ 27{,}21\)&\(5\ 7{,}64\)&\(120^0 51\ 21{,}48\)&\(54\ 3{,}80\)&\(58\ 36{,}95\)
\end{tabular}
\endgroup
\end{center}
und hieraus nach Abzug der obigen Störungen \(\delta\AscendingNode\), \(\delta i\) etc. die mittlern Elemente:
\begin{center}
\begingroup\fontsize{8}{10}\selectfont\setlength{\tabcolsep}{3pt}
\begin{tabular}{@{}lr|r|r|r|r|r|r|r@{}}
\(\AscendingNode=\)&\(172^0 31'\,45\overset{\prime\prime}{,}91\)&\(45\overset{\prime\prime}{,}53\)&\(46\overset{\prime\prime}{,}60\)&\(46\overset{\prime\prime}{,}37\)&\(44\overset{\prime\prime}{,}70\)&\(43\overset{\prime\prime}{,}32\)&\(43\overset{\prime\prime}{,}63\)&\(43\overset{\prime\prime}{,}97\)\\
\(i=\)&\(34\ 36\ 9{,}74\)&\(10{,}23\)&\(10{,}25\)&\(9{,}20\)&\(9{,}70\)&\(11{,}50\)&\(12{,}22\)&\(10{,}44\)\\
\(\varphi=\)&\(14\ 2\ 55{,}08\)&\(54{,}97\)&\(54{,}93\)&\(54{,}40\)&\(52{,}70\)&\(46{,}76\)&\(46{,}83\)&\(46{,}44\)\\
\(\tilde\omega=\)&\(121\ 7\ 42{,}76\)&\(44{,}60\)&\(49{,}49\)&\(55{,}81\)&\(61{,}14\)&\(56{,}67\)&\(68{,}65\)&\(73{,}73.\)
\end{tabular}
\endgroup
\end{center}
Im Mittel also:
\begin{center}
\begin{tabular}{@{}lcrlrr@{}}
\(\AscendingNode\) medius&1810&\(:\)&\(172^0 31'\,45\overset{\prime\prime}{,}00\)&mittlerer Fehler:&\(1\overset{\prime\prime}{,}20\)\\
Incl. media&»&\(:\)&\(34\quad36\quad10{,}41\)&»\qquad »&\(0{,}93\)\\
\(\varphi\)&»&\(:\)&\(14\quad2\quad51{,}51\)&»\qquad »&\(3{,}81\)\\
\(\tilde\omega\)&»&\(:\)&\(121\quad7\quad56{,}61.\)&&
\end{tabular}
\end{center}
\SourceLeaf{567}
Diese Werthe zeigen, wie über Erwarten scharf die Übereinstimmung zwischen den Resultaten der speciellen und der generellen Störungen ist.]
\endgroup
\clearpage
\endgroup
\end{landscape}
